\section{Introduction}
\label{sec:intro}

The number of bits per weight is the only knob that changes the \emph{class}
of model a given machine can load. At 16 bits a 70B-parameter model needs
140~GB; at 2 bits it needs ${\sim}20$~GB and fits in the unified memory of a
workstation. For anyone whose constraint is sovereignty (capable models on hardware
they own), the 2-bit operating point determines whether the model loads at
all.

The best published quality at 2 bits comes from Leech-lattice vector
quantization (LLVQ)~\citep{llvq2026}: weights are grouped into blocks of 24,
normalized, and snapped to the nearest point of the Leech lattice
$\Lambda_{24}$, whose exceptional packing density buys roughly a decibel of
distortion over trellis and E8-based codes at the same rate. But the paper's
own systems story stops short. Its CUDA kernel decodes a single lattice shell
($m=3$, ``for simplicity''), is reported slower than QTIP's
kernel~\citep{qtip2024}, and the authors explicitly defer low-level
optimization as ``largely orthogonal'' to their contribution. The decoder the
2-bit rate actually requires is a union of lattice shells with hundreds of
distinct coordinate-permutation classes (301 for the paper's best
configuration, the ball $\Lambda_{24}(12)$). It exists nowhere, including
in the original work.

We build it, and then test the claim that motivated it. Ported into our
harness at a pinned commit and timed in the same rounds as everything else,
QTIP's kernel runs $2.27\times$ faster than our best layout: low-level
optimization is not ``largely orthogonal'' to this contribution, it
\emph{is} the contribution, and on that axis the lattice loses. The
mechanism is neither rate nor quality --- the two formats match on both ---
but the load-time unfolding a combinatorial lattice index demands and a
sliding-window trellis does not (\S\ref{sec:qtip}). Reaching that sentence
took a complete implementation of both sides.

This paper is about that decoder, and about what it takes to move
lattice-decoded weights from merely smaller to \emph{faster} than FP16.
Three problems have to be solved together:

\begin{enumerate}
\item \textbf{The decode itself} (\S\ref{sec:decoder}). A $\Lambda_{24}$
index is 47 bits naming one of $1.1\times10^{14}$ lattice points. Decoding it
inside a matvec inner loop must cost shifts and masks, not arithmetic, and
must not diverge across a warp whose 32 lanes hold different classes.

\item \textbf{The memory format} (\S\ref{sec:layouts}). The rate that
determines what fits on a GPU is the rate of whatever format the kernel reads,
a different quantity from the on-disk rate (2.07 bits per weight
of payload; \S\ref{sec:evaluation}). Our
first working layout cost 5.51 bits per weight in VRAM, \emph{more} than
a 4-bit format, erasing the point of 2-bit quantization. We show a
progression of layouts that brings the in-VRAM rate to 4.80 and then 4.34
bits per weight at bit-exact decoded content while \emph{increasing}
throughput, and one dead end (3.59, compute-bound) that maps the floor. We
attribute what the kernel actually waits on with an instrumented job, and
place two deployed competitors on the same axis in the same process: a
4-bit kernel that converts 88\% of its byte advantage into time against our
65\%, and the 2-bit kernel above. The latter also runs faster than a
control of ours that reads \emph{no weight bytes at all}, which retires
that control as a floor: what we had measured as a ceiling on any format
turns out to be the cost of our own launch geometry.

\item \textbf{The integration} (\S\ref{sec:integration}). End-to-end
throughput is gated by the serving engine, not the kernel. We found and
worked around a per-token 778~MB transposed materialization in the output
head of a widely used Rust inference engine, and we report our speedup in
a double formulation
that separates the kernel's contribution from the engine fix.
\end{enumerate}

The result (\S\ref{sec:evaluation}): a 2-bit Qwen3-4B served at 87.0~tok/s
in 2.60~GB against 43.5~tok/s in 8.04~GB for the FP16 reference, with the
serving stack costing nothing measurable on top of the quantized artifact.
The same path serves a 2-bit Qwen3-14B at 43.3~tok/s in 9.39~GB against
17.0~tok/s in 29.54~GB, and then stops one model class later --- the
activation rotation is a Walsh--Hadamard transform, so a single block must
hold the whole activation in shared memory, and the 32B exceeds the
measured 101\,376-byte ceiling by 1\,024 bytes (\S\ref{sec:limitations}).
We are equally explicit about what this does not buy
(\S\ref{sec:limitations}).
At 4B, 2 bits costs $\times1.384$ perplexity and 14.6 MMLU points (served
arm; 14.73 on the sealed artifact), a deficit
concentrated in reasoning-heavy subjects, while the 4-bit baseline keeps
essentially full quality: on that model, 4 bits wins. The case for 2 bits
rests on scale, and we measure it on two further models. The degradation
falls ($\times1.384 \to \times1.220 \to \times1.189$ perplexity across 4B, 8B
and 14B, the MMLU gap to 4 bits going 14.45, 7.49, 6.09 points, all paired
estimates with 95\% intervals), and it flattens as it falls: the excess
over FP16 drops by 42.8\% on the first step and 13.9\% on the
second. Whether that flattening is \emph{established} depends on the metric:
resolved on perplexity, where the evaluation windows pair across sizes;
unresolved on the MMLU gap, likely a matter of statistical power. We
carry the two verdicts separately throughout
(\S\ref{sec:evaluation}, Appendix~\ref{app:scale}). On memory the three
points are consistent: the served model is below the 4-bit
baseline at every size, by a margin that peaks in the middle because it
tracks the embedding share.

Every figure regenerates from CSV files committed in the repository, the
layout table is checked against its CSV on every build, and each number
traces to a dated GPU job with its billed cost.
